% GENERATED FILE. Edit canonical lessons, then run npm run generate:books.
% canonical-source-hash: fnv1a64:833cebf959910472
% canonical-lessons: UR-C279-pahchanna, UR-C279-huliya, UR-C279-shubah, UR-C279-manzar, UR-C279-nazar-ana

\chapter{What the Witness Saw}
\label{ch:ur-what-the-witness-saw}
\hlchaptermodality{279}

\begin{chapteropening}
\textbf{By the end of this chapter:} \emph{I can say what I saw and what would have been different.}

Complete the last lesson of chapter 279: I can say what I saw and what would have been different.
\end{chapteropening}

\section[\texorpdfstring{\ur{پہچاننا} (pahchānnā)}{پہچاننا (pahchānnā)}]{\ur{پہچاننا} (pahchānnā) — to recognize}

\label{lesson:UR-C279-pahchanna}



\begin{quote}
Before the new one: say the Urdu for a lock-up, then the Urdu for a statement.
\end{quote}

\subsection*{You'll want to know: \ur{پہچاننا}}
\textbf{\ur{پہچاننا}} — \emph{pahchānnā} — \textquotedblleft{}to recognize\textquotedblright{}.

\begin{quote}
\ur{اگر} \ur{میں} \ur{قریب} \ur{ہوتا} \ur{تو} \ur{اسے} \ur{پہچانتا۔}
\end{quote}

\begin{quote}
\emph{agar maiṅ qarīb hotā to use pahchāntā.}
\end{quote}

\begin{quote}
If I had been closer, I would have recognized him.
\end{quote}

Past, with \ur{نے} (he or she …): \ur{اس} \ur{نے} \ur{چور} \ur{کو} \ur{پہچانا}, \emph{us ne chor ko pahchānā}.

\begin{cousinweb}
\textbf{\ur{پہچاننا}} (\emph{pahchānnā}) is inherited: CDIAL lists the Hindi cousin \emph{pahcānnā} under Sanskrit \emph{pratyabhijānāti}, \textquotedblleft{}recognizes\textquotedblright{} (CDIAL 8637). A person recognized takes \textbf{\ur{کو}}, and the verb stays masculine: \textbf{\ur{چور} \ur{کو} \ur{پہچانا}}.
\end{cousinweb}

\begin{grammarlens}[title={What you've built}]
One more word for telling what happened.
\end{grammarlens}

\subsection*{Your turn}
\begin{itemize}
\raggedright
  \item \emph{Say it:} \emph{pahchānnā}
  \item \emph{Say it:} \emph{pahchānnā}, once more
  \item \emph{Say it:} \emph{agar maiṅ qarīb hotā to use pahchāntā.} — If I had been closer, I would have recognized him.
  \item \emph{Recall:} say the Urdu for a lock-up, then the Urdu for a statement, then say \emph{pahchānnā} again
\end{itemize}

\begin{tcolorbox}[breakable,colback=teal!4,colframe=teal!35!black,title={Before you move on}]
What does \ur{پہچاننا} mean? (\textquotedblleft{}To recognize\textquotedblright{}.) Say it once more.
\end{tcolorbox}
\section[\texorpdfstring{\ur{حلیہ} (huliyā)}{حلیہ (huliyā)}]{\ur{حلیہ} (huliyā) — a person's description; someone's looks}

\label{lesson:UR-C279-huliya}



\begin{quote}
Before the new one: say the Urdu for a statement, then the Urdu for to recognize.
\end{quote}

\subsection*{You'll want to know: \ur{حلیہ}}
\textbf{\ur{حلیہ}} — \emph{huliyā} — \textquotedblleft{}a person's description; someone's looks\textquotedblright{}.

\begin{quote}
\ur{گواہ} \ur{نے} \ur{چور} \ur{کا} \ur{حلیہ} \ur{بتایا۔}
\end{quote}

\begin{quote}
\emph{gavāh ne chor kā huliyā batāyā.}
\end{quote}

\begin{quote}
The witness described what the thief looked like.
\end{quote}

Masculine.

\begin{cousinweb}
is usually traced, through Persian, to Arabic \emph{ḥilya}, \textquotedblleft{}an adornment, the look of a person\textquotedblright{}. The police ask for the \textbf{\ur{حلیہ}}: height, face, clothes.
\end{cousinweb}

\begin{grammarlens}[title={What you've built}]
One more word for telling what happened.
\end{grammarlens}

\subsection*{Your turn}
\begin{itemize}
\raggedright
  \item \emph{Say it:} \emph{huliyā}
  \item \emph{Say it:} \emph{huliyā}, once more
  \item \emph{Say it:} \emph{gavāh ne chor kā huliyā batāyā.} — The witness described what the thief looked like.
  \item \emph{Recall:} say the Urdu for a statement, then the Urdu for to recognize, then say \emph{huliyā} again
\end{itemize}

\begin{tcolorbox}[breakable,colback=teal!4,colframe=teal!35!black,title={Before you move on}]
What does \ur{حلیہ} mean? (\textquotedblleft{}A person's description; someone's looks\textquotedblright{}.) Say it once more.
\end{tcolorbox}
\section[\texorpdfstring{\ur{شبہ} (shubah)}{شبہ (shubah)}]{\ur{شبہ} (shubah) — a suspicion}

\label{lesson:UR-C279-shubah}



\begin{quote}
Before the new one: say the Urdu for to recognize, then the Urdu for a person's description.
\end{quote}

\subsection*{You'll want to know: \ur{شبہ}}
\textbf{\ur{شبہ}} — \emph{shubah} — \textquotedblleft{}a suspicion\textquotedblright{}.

\begin{quote}
\ur{مجھے} \ur{پڑوسی} \ur{پر} \ur{شبہ} \ur{ہے۔}
\end{quote}

\begin{quote}
\emph{mujhe paṛosī par shubah hai.}
\end{quote}

\begin{quote}
I suspect the neighbour.
\end{quote}

Masculine.

\begin{cousinweb}
\textbf{\ur{شبہ}} (\emph{shubah}) is usually traced, through Persian, to Arabic \emph{shubha}, \textquotedblleft{}a doubt, an uncertainty\textquotedblright{}. The person suspected takes \textbf{\ur{پر}}: \textbf{\ur{پڑوسی} \ur{پر} \ur{شبہ}}.
\end{cousinweb}

\begin{grammarlens}[title={What you've built}]
One more word for telling what happened.
\end{grammarlens}

\subsection*{Your turn}
\begin{itemize}
\raggedright
  \item \emph{Say it:} \emph{shubah}
  \item \emph{Say it:} \emph{shubah}, once more
  \item \emph{Say it:} \emph{mujhe paṛosī par shubah hai.} — I suspect the neighbour.
  \item \emph{Recall:} say the Urdu for to recognize, then the Urdu for a person's description, then say \emph{shubah} again
\end{itemize}

\begin{tcolorbox}[breakable,colback=teal!4,colframe=teal!35!black,title={Before you move on}]
What does \ur{شبہ} mean? (\textquotedblleft{}A suspicion\textquotedblright{}.) Say it once more.
\end{tcolorbox}
\section[\texorpdfstring{\ur{منظر} (manzar)}{منظر (manzar)}]{\ur{منظر} (manzar) — a scene; a sight}

\label{lesson:UR-C279-manzar}



\begin{quote}
Before the new one: say the Urdu for a person's description, then the Urdu for a suspicion.
\end{quote}

\subsection*{You'll want to know: \ur{منظر}}
\textbf{\ur{منظر}} — \emph{manzar} — \textquotedblleft{}a scene; a sight\textquotedblright{}.

\begin{quote}
\ur{میں} \ur{وہ} \ur{منظر} \ur{نہیں} \ur{بھولوں} \ur{گا۔}
\end{quote}

\begin{quote}
\emph{maiṅ voh manzar nahīṅ bhūlūṅ gā.}
\end{quote}

\begin{quote}
I will not forget that scene.
\end{quote}

Masculine.

\begin{cousinweb}
\textbf{\ur{منظر}} (\emph{manzar}) is usually traced, through Persian, to Arabic \emph{manẓar}, \textquotedblleft{}a place looked at, a sight\textquotedblright{}. It is what lay before the eyes at one moment.
\end{cousinweb}

\begin{grammarlens}[title={What you've built}]
One more word for telling what happened.
\end{grammarlens}

\subsection*{Your turn}
\begin{itemize}
\raggedright
  \item \emph{Say it:} \emph{manzar}
  \item \emph{Say it:} \emph{manzar}, once more
  \item \emph{Say it:} \emph{maiṅ voh manzar nahīṅ bhūlūṅ gā.} — I will not forget that scene.
  \item \emph{Recall:} say the Urdu for a person's description, then the Urdu for a suspicion, then say \emph{manzar} again
\end{itemize}

\begin{tcolorbox}[breakable,colback=teal!4,colframe=teal!35!black,title={Before you move on}]
What does \ur{منظر} mean? (\textquotedblleft{}A scene; a sight\textquotedblright{}.) Say it once more.
\end{tcolorbox}
\section[\texorpdfstring{\ur{نظر} \ur{آنا} (nazar ānā)}{نظر آنا (nazar ānā)}]{\ur{نظر} \ur{آنا} (nazar ānā) — to be seen; to come into view}

\label{lesson:UR-C279-nazar-ana}



\begin{quote}
Before the new one: say the Urdu for a suspicion, then the Urdu for a scene.
\end{quote}

\subsection*{You'll want to know: \ur{نظر} \ur{آنا}}
\textbf{\ur{نظر} \ur{آنا}} — \emph{nazar ānā} — \textquotedblleft{}to be seen; to come into view\textquotedblright{}.

\begin{quote}
\ur{مجھے} \ur{کھڑکی} \ur{میں} \ur{ایک} \ur{آدمی} \ur{نظر} \ur{آیا۔}
\end{quote}

\begin{quote}
\emph{mujhe khiṛkī meṅ ek ādmī nazar āyā.}
\end{quote}

\begin{quote}
I saw a man at the window.
\end{quote}

Past (he …): \ur{وہ} \ur{نظر} \ur{آیا}, \emph{voh nazar āyā}.

\begin{cousinweb}
\textbf{\ur{نظر} \ur{آنا}} (\emph{nazar ānā}) is built from \emph{nazar}, a look, usually traced, through Persian, to Arabic \emph{naẓar}, \textquotedblleft{}a looking\textquotedblright{}, and \textbf{\ur{آنا}}: a thing comes into your look. The one who sees takes \textbf{\ur{مجھے}}, and the verb follows the thing seen: \textbf{\ur{آدمی} \ur{نظر} \ur{آیا}}.
\end{cousinweb}

\begin{grammarlens}[title={What you've built}]
One more word for telling what happened.
\end{grammarlens}

\subsection*{Your turn}
\begin{itemize}
\raggedright
  \item \emph{Say it:} \emph{nazar ānā}
  \item \emph{Say it:} \emph{nazar ānā}, once more
  \item \emph{Say it:} \emph{mujhe khiṛkī meṅ ek ādmī nazar āyā.} — I saw a man at the window.
  \item \emph{Recall:} say the Urdu for a suspicion, then the Urdu for a scene, then say \emph{nazar ānā} again
\end{itemize}

\begin{tcolorbox}[breakable,colback=teal!4,colframe=teal!35!black,title={Before you move on}]
What does \ur{نظر} \ur{آنا} mean? (\textquotedblleft{}To be seen; to come into view\textquotedblright{}.) Say it once more.
\end{tcolorbox}
